% Содержимое отчета по курсу Анализ алгоритмов

\aaunnumberedsection{ВВЕДЕНИЕ}{sec:intro}

Параллельные вычисления по конвейерному принципу представляют собой подход,
при котором обработка задачи разбивается на несколько независимых этапов (стадий),
выполняемых последовательно, но параллельно. Каждая стадия обрабатывает свою часть данных,
передавая результаты на следующую стадию через общие очереди.
Это позволяет эффективно использовать многопоточность,
минимизируя простой процессоров за счёт равномерного распределения
нагрузки между потоками~\cite{paralel_cpp}.

Нативные потоки (Native Threads) — это потоки, управляемые непосредственно операционной системой.
Использование нативных потоков позволяет создавать высокопроизводительные многопоточные приложения,
но требует контроля доступа к общим ресурсам для предотвращения гонок данных.

Цель работы --- получение навыка организации параллельных вычислений по конвейерному принципу.

Для достижения поставленной цели необходимо выполнить следующие задачи:
\begin{itemize}
    \item описать входные и выходные данные;
    \item анализ структуры сайта;
    \item разработка ПО, выполняющего парсинг рецептов в конвейере и загрузку информации в базу данных;
    \item исследование характеристик разработанного ПО на данных о интервалах времени стадий обработки рецепта.
\end{itemize}

\aasection{Входные и выходные данные}{sec:input-output}

Входными данными программы являются скаченные страницы рецептов с ресурса~\cite{recepies}.
Выходными данными является файлы json -- загруженные рецепты по ссылкам из входных данных, при этом в каждом json файле хранится:

\begin{itemize}
    \item id задачи;
    \item название рецепта;
    \item список ингредиентов;
    \item шаги рецепта.
\end{itemize}

\aasection{Преобразование входных данных в выходные}{sec:algorithm}

Входные данные, представленные в виде текста страницы рецепта, проходят три этапа обработки.
На первом этапе выполняется чтение страницы из файла, где извлекается её содержимое.
На втором этапе текст анализируется: выделяются URL рецепта, название, ингредиенты
(с указанием названия, количества и единицы измерения), шаги приготовления и URL изображения.
Эти данные очищаются от HTML-тегов и других лишних символов.
На третьем этапе обработанная информация преобразуется в формат,
пригодный для хранения в базе данных (JSON для ингредиентов и шагов) и
записывается в соответствующие поля таблицы.



\aasection{Тестирование}{sec:tests}

Тестирование программы проводилось подачей на вход директории с 60 файлами со страницами различных
рецептов веб-ресурса~\cite{recepies}. При этом отслеживалось количество
удачных скачиваний и обработок страницы, а также выборочная ручная проверка 5 случайных выходных файлов.

\aasection{Описание исследования}{sec:study}

В ходе исследования требуется сформировать лог обработки задач.


\section{Технические характеристики}
Технические характеристики используемого устройства:
\begin{itemize}
    \item процессор: AMD Ryzen™ 7--5800H with Radeon™ Graphics × 16;
    \item ОС: 64-разрядная операционная система Ubuntu 24.04.1 LTS;
    \item оперативная память: 16,0 ГБ;
    \item число логических ядер -- 16.
\end{itemize}


\section{Результаты исследования}

В таблице~\ref{tbl:b_log} приведен фрагмент лога обработки. Обозначения событий:
\begin{itemize}
    \item start\_step\_1 — начло обработки на стадии 1;
    \item end\_step\_1 — окончание обработки на стадии 1;
    \item start\_step\_2 — начло обработки на стадии 2;
    \item end\_step\_2 — окончание обработки на стадии 2;
    \item start\_step\_3 — начло обработки на стадии 3;
    \item end\_step\_3 — окончание обработки на стадии 3.
\end{itemize}

\begin{longtable}{|p{.33\textwidth - 2\tabcolsep}|p{.33\textwidth - 2\tabcolsep}|p{.34\textwidth - 2\tabcolsep}|}
    \caption{Фрагмент лога обработки (начало)}\label{tbl:b_log}
    \\
    \hline
    Метка времени, мкс & Событие        & ID записи \\
    \hline
    \endfirsthead
    \caption{Фрагмент лога обработки (окончание)}
    \\
    \hline
    Метка времени, мкс   & Событие   & ID записи        \\
    \hline
    \endhead
    \hline
    \endfoot
    \endlastfoot
    \hline
    1732747886220541   & start\_step\_1 & 32        \\\hline
    1732747886220575   & end\_step\_1   & 32        \\\hline
    1732747886220593   & start\_step\_1 & 33        \\\hline
    1732747886220628   & end\_step\_1   & 33        \\\hline
    1732747886220644   & start\_step\_1 & 34        \\\hline
    1732747886220680   & end\_step\_1   & 34        \\\hline
    1732747886220700   & start\_step\_1 & 35        \\\hline
    1732747886220732   & end\_step\_1   & 35        \\\hline
    1732747886220752   & start\_step\_1 & 36        \\\hline
    1732747886220785   & end\_step\_1   & 36        \\\hline
    1732747886220806   & start\_step\_1 & 37        \\\hline
    1732747886220842   & end\_step\_1   & 37        \\\hline
    1732747886220862   & start\_step\_1 & 38        \\\hline
    1732747886220898   & end\_step\_1   & 38        \\\hline
    1732747886220920   & start\_step\_1   & 39         \\\hline
    1732747886220954   & end\_step\_1 & 39         \\\hline
    1732747886240923   & end\_step\_2 & 1         \\\hline
    1732747886240957   & start\_step\_2   & 2         \\\hline
    1732747886241016   & start\_step\_3   & 1         \\\hline
    1732747886243748   & end\_step\_3   & 1         \\\hline
\end{longtable}

В ходе исследования получены следующие характеристики:
\begin{itemize}
    \item среднее время жизни задачи -- 448332 мкс;
    \item среднее время выполнения задачи 1 -- 49 мкс;
    \item среднее время выполнения задачи 2 -- 21540 мкс;
    \item среднее время выполнения задачи 3 -- 1959 мкс;
    \item среднее время ожидания в очереди после шага 1 -- 423000 мкс;
    \item среднее время ожидания в очереди после шага 2 -- 51 мкс;
    \item среднее время ожидания в очереди после шага 3 -- 49 мкс.
\end{itemize}

При конвейерной многопоточной обработке сложность выполнения задачи и время ожидания
в очереди тесно взаимосвязаны. Увеличение сложности задачи,
требующей больше времени на выполнение, неизбежно приводит к увеличению
времени ожидания последующих задач в очередях. Это связано с тем, что более
сложные этапы обработки становятся узкими местами в конвейере,
замедляя прохождение задач через систему.

\clearpage

\aaunnumberedsection{ЗАКЛЮЧЕНИЕ}{sec:end}

Многопоточная конвейерная обработка является эффективным методом
организации параллельных вычислений, особенно для задач,
требующих последовательной обработки данных через несколько этапов.
Основное преимущество этого подхода заключается в том, что каждую стадию
обработки можно выполнять независимо в отдельных потоках, что позволяет
повысить производительность за счёт параллельной работы.

Цель работы была достигнута.
Решены все поставленные задачи:
\begin{itemize}
    \item описаны входные и выходные данные;
    \item разработано ПО, выполняющее парсинг рецептов в конвейере и загрузку информации в базу данных;
    \item выполнено исследование характеристик разработанного ПО на данных о интервалах времени стадий обработки рецепта.
\end{itemize}
